\subsection{Равновесие твердого тела}
%%%%%%%%%%%%%%%%%%%%%%%%%%%%%%%%%%%%%%%%%%%%%%%%%%%%
%%%%%%%%%%%%%%%%%%%%%%%%%%%%%%%%%%%%%%%%%%%%%%%%%%%%
 \z[{\begin{wrapfigure}[3]{r}{0.2\textwidth}%
		\vspace{-12pt}%
		\begin{flushright}%
 			\includegraphics[width=0.2\textwidth]{8.2.1}%
		\end{flushright}%
	\end{wrapfigure}%
	}]%
	{Каким должен быть коэффициент трения $\mu$ однородного стержня о пол, чтобы он мог стоять так, как показано на рисунке? Длина нити $AB$ равна длине стержня.
\vspace{3em}
	}
%
{$\mu = 0.33$}
%%%%%%%%%%%%%%%%%%%%%%%%%%%%%%%%%%%%%%%%%%%%%%%%%%%%
\z[{\begin{wrapfigure}[3]{r}{0.2\textwidth}%
		\vspace{-10pt}%
		\begin{flushright}%
 			\includegraphics[width=0.2\textwidth]{8.2.2}%
		\end{flushright}%
	\end{wrapfigure}%
	}]%
	{С какой силой $N$ давит на стенки цилиндрического стакана палочка массы $m$, наполовину погруженная в воду? Угол наклона палочки к горизонтали $\alpha$. Трением пренебречь.
\vspace{3em}
	}
%
{$F = \frac{1}{4}mg \ctg \alpha$}
%%%%%%%%%%%%%%%%%%%%%%%%%%%%%%%%%%%%%%%%%%%%%%%%%%%% 
 \z[{\begin{wrapfigure}[5]{r}{0.3\textwidth}%
		\vspace{-15pt}%
		\begin{flushright}%
 			\includegraphics[width=0.3\textwidth]{8.2.3}%
		\end{flushright}%
	\end{wrapfigure}%
	}]%
	{Расстояние между вертикальными стенками равно $l$. Какой длины $L$ стержень, вставленный наискось между стенками, не будет опускаться, если коэффициент трения между стержнем и стенками равен $\mu$?
\vspace{3em}
	}
%
{$l < L < l \sqrt{1 + \mu^2}$}
%%%%%%%%%%%%%%%%%%%%%%%%%%%%%%%%%%%%%%%%%%%%%%%%%%%% 
 \z[{\begin{wrapfigure}[4]{r}{0.2\textwidth}%
		\vspace{-10pt}%
		\begin{flushright}%
 			\includegraphics[width=0.2\textwidth]{8.2.4}%
		\end{flushright}%
	\end{wrapfigure}%
	}]%
	{Однородный куб с помощью веревки, привязанной в середине его ребра, подвешен к вертикальной стене. При каких значениях угла $\varphi$ между веревкой и стенкой куб соприкасается со стенкой всей гранью, если коэффициент трения его о плоскость равен $\mu$?
\vspace{3em}
	}
%
{$\tan \alpha \leq \frac{1}{\mu}$}
%%%%%%%%%%%%%%%%%%%%%%%%%%%%%%%%%%%%%%%%%%%%%%%%%%%%
%%%%%%%%%%%%%%%%%%%%%%%%%%%%%%%%%%%%%%%%%%%%%%%%%%%%	
 \z[{\begin{wrapfigure}[3]{r}{0.2\textwidth}%
		\vspace{-10pt}%
		\begin{flushright}%
 			\includegraphics[width=0.2\textwidth]{8.2.5}%
		\end{flushright}%
	\end{wrapfigure}%
	}]%
	{С какой минимальной силой $F$ необходимо тянуть колесо радиуса $R$ за ось вращения, чтобы поднять его на ступеньку высотой $h$ ($R > h$)?}
 %
{$F \geq mg \sqrt{\frac{h(2R - h)}{(R - h)^2}}$}